%! Comparing Distributions — Unit 2, Lesson 2.4 — Guided Notes & Practice
% THE in-class document, in the MAIN IDEAS / NOTES shape (LESSON_SHAPE.md §1): the vocabulary
% box, the hook, then ONE two-column guidednotes table. Each row is one idea — a label on the
% left; on the right one or two COMPLETE printed sentences, ONE large pre-drawn display the
% student reads or names, then a bold prompt and a two-across grid of numbered problems.
% Problems run 1..21. Unit 2 timing (5/45/5/5): I do ~25 min (rows 1-4), Guided Practice ~20.
% DENSITY: a \blank{} only where a word or number IS the answer (the eight cells of the
% comparison table, problem 9 — the one \wordbank strip sits above it); no mid-sentence blanks;
% every display is pre-drawn; the page belongs to the student's pen.
%   I DO   : the teacher reads the definition, marks up the display, works the first problem
%            of each grid (1, 5, 10, 13).
%   WE DO  : the class works the rest of each grid, students holding the pen; the last row,
%            Guided Practice (17-20), is worked entirely together on paper airplanes.
%   YOU DO : nothing on this page — ../homework/ IS the individual practice, assigned at Close
%            & Assign and done outside class.
% THE TRAPS: problem 3 (a left-side leaf read backwards, 81 for 18) and problem 8 (two boxplots
% on different scales — the secondary misconception).
% THE CRUX: problem 14 (row 4, "Different Spread"): two shops with the SAME median, and the
% hook vote re-taken — spread decides. Problem 16 is the non-comparison (two descriptions with
% no comparative words). Problem 19 is the crux transferred (same mean, different SD).
%
% THE DATA (quartiles: median of each half, overall median left out when n is odd)
%   Tony's Pizza (n=11): 18 21 22 24 25 25 27 28 30 33 44
%     min 18, Q1 22, med 25, Q3 30, max 44; IQR 8; fences 10 and 42 -> 44 is an outlier;
%     mean 297/11 = 27.0 (> median: skewed right)
%   Rosa's Pizza (n=11): 24 29 31 33 34 35 36 36 37 38 39
%     min 24, Q1 31, med 35, Q3 37, max 39; IQR 6; fences 22 and 46 -> no outliers;
%     mean 372/11 = 33.8 (< median: skewed left)
%   Slice House (n=11): 15 20 24 28 33 35 38 44 47 50 55
%     min 15, Q1 24, med 35, Q3 47, max 55; IQR 23; fences -10.5 and 81.5 -> no outliers
%     within 40 min: Rosa's 11/11, Slice House 7/11; under 25 min: Rosa's 1/11, Slice 3/11
%   Guided Practice, paper airplanes (n=10 each, feet):
%     Design A: 26 27 28 29 30 30 31 32 33 34 -> mean 30.0, Sx = sqrt(60/9) = 2.58;
%       Q1 28, med 30, Q3 32; in the 27-33 zone: 8 of 10
%     Design B: 18 22 25 28 30 30 32 35 38 42 -> mean 30.0, Sx = sqrt(474/9) = 7.26;
%       Q1 25, med 30, Q3 35; in the 27-33 zone: 4 of 10
%
% PAGE LOCKSTEP with notes_key/: the key differs only in -key for -boxes, the header's
% "--- Answer Key", \vterm -> \vtermans, \blank -> \ans, \writelines{2} -> two \ansline, and
% the answer argument of every \pcell, \writespace and \labelbox. KEEP EVERY \pcell ON ONE
% SOURCE LINE.
% TABLE MECHANICS: inside a cell, `\\` ENDS THE ROW — break lines with \par. A row cannot break
% across pages: the figure gets its own sub-row (`\\` then `& ...`) and EACH GRID ROW is its own
% sub-row — close the probgrid, `\\ &`, reopen as probgrid* (no top rule).
% NO SKETCH QUESTIONS: every display is pre-drawn in TikZ (or set as a table) and read.
\documentclass[12pt]{article}
\usepackage{saar-article}
\usepackage{saar-boxes}
\newcommand{\TallMath}[1]{$\displaystyle #1\rule[-1.4em]{0pt}{3.2em}$}
% Word bank strip — ONLY above a table the student fills (row 3). Defined per document;
% the key inherits it and prints the same strip, so heights match.
\newcommand{\wordbank}[1]{\par\vspace{2pt}\noindent\fcolorbox{goldacc}{goldbg}{\parbox{\dimexpr\linewidth-2\fboxsep-2\fboxrule\relax}{\small\textbf{Word bank:} #1}}\par\vspace{4pt}}
% Vocabulary rows: a FIXED-HEIGHT row carrying the term label and OPEN writing space —
% no underline beside the term and no rule line (user correction 2026-09-06). The key fills
% exactly the same height, so the box cannot drift blank vs. keyed. Copied from
% unit01/lesson02/notes/main.tex. Keep key definitions to two lines.
\newcommand{\termrowheight}{2.3cm}
\newlength{\termrowinset}\setlength{\termrowinset}{7pt}
\newcommand{\vterm}[1]{%
  \par\noindent\begin{minipage}[t][\termrowheight][t]{\linewidth}%
    \vspace*{\termrowinset}%
    \textbf{\textcolor{cerulean}{#1:}}%
  \end{minipage}\par}
\newcommand{\vtermans}[2]{%
  \par\noindent\begin{minipage}[t][\termrowheight][t]{\linewidth}%
    \vspace*{\termrowinset}%
    \textbf{\textcolor{cerulean}{#1:}}\quad\ans{#2}%
  \end{minipage}\par}

\begin{document}

\pageheader{Unit 2, Lesson 2.4}{Guided Notes \& Practice}
% NAMESTRIP: no \namedateperiod here. The student writes their name once, on the
% cover this component is stapled behind.

\vspace{0.15in}

\boxguard[16]
\begin{vocabbox}
\small
Fill in each term \textbf{as we name it} in the notes below.
\vterm{Back-to-back stemplot}
\vterm{Parallel boxplots}
\vterm{Common scale}
\vterm{Comparison}
\vterm{Consistent}
\end{vocabbox}

\boxguard[12]
\begin{hookbox}
\small
Rosa's Pizza and Slice House advertise the same thing: \emph{a median delivery time of
$35$ minutes.} Your party starts in $40$ minutes, and the pizza has to be there on time.

\vspace{3pt}
\textbf{Which shop should you call?} Circle one: \quad Rosa's \qquad Slice House \qquad
it doesn't matter --- same median \quad --- and say why. We vote again at problem~14.
\par\writelines{2}
\end{hookbox}

\small
\begin{guidednotes}
% ── ROW 1: back-to-back stemplots ────────────────────────────────────────────
\mainidea[Two groups, one stem]{Back-to-Back Stemplots} &
A \textbf{back-to-back stemplot} puts two groups on one shared column of stems: one group's
leaves grow to the left, the other group's to the right. \textbf{Read every leaf with its
stem} --- on either side, the stem is the tens digit and the leaf is the ones digit.
\\
& {\centering
{\footnotesize\textbf{Delivery times (minutes), $11$ orders from each shop}}\par\vspace{4pt}
{\large\ttfamily\renewcommand{\arraystretch}{1.3}%
\begin{tabular}{r|c|l}
\multicolumn{1}{c}{\normalsize\rmfamily\bfseries Tony's Pizza} & \multicolumn{1}{c}{\normalsize\rmfamily\bfseries Stem} & \multicolumn{1}{c}{\normalsize\rmfamily\bfseries Rosa's Pizza} \\ \hline
8 & 1 & \\
8 7 5 5 4 2 1 & 2 & 4 9 \\
3 0 & 3 & 1 3 4 5 6 6 7 8 9 \\
4 & 4 & \\ \hline
\end{tabular}}\par\vspace{6pt}
Key: $5 \,|\, 2 \,|\, 4$ means Tony's \labelbox{1.8cm}{}\quad and Rosa's \labelbox{1.8cm}{}\par}
\\
& \notesprompt{Read both sides --- stem first, then leaf.}
\begin{probgrid}{|Y|Y|}
\pcell{1}{Read each shop's fastest and slowest delivery.}{1.6cm}{} & \pcell{2}{Find each shop's median delivery time. (Eleven orders: count to the $6$th leaf from the fastest.)}{1.6cm}{} \\ \hline
\end{probgrid}
\\
& \begin{probgrid}*{|Y|Y|}
\pcell{3}{Jaden reads the top leaf on Tony's side as $81$ minutes. What did he do wrong, and what is the value?}{2.0cm}{} & \pcell{4}{Describe each shop's shape. Does either shop have a value that stands apart?}{2.0cm}{} \\ \hline
\end{probgrid}
\\ \hline
% ── ROW 2: parallel dot plots and parallel boxplots on a common scale ────────
\mainidea[One number line]{Parallel Displays} &
\textbf{Parallel dot plots} and \textbf{parallel boxplots} stack the groups over \textbf{one
common scale}, so the same distance means the same number of minutes for every group.
Quartiles are found as in Lesson~2.3: the median of each half, leaving out the overall median
when $n$ is odd.
\\
& {\centering
\begin{tikzpicture}[x=0.3cm, y=1cm]
  % one common scale: faint gridlines run through every display
  \foreach \v in {15,20,...,45} { \draw[linegray, densely dashed] (\v,0) -- (\v,4.35); }
  \draw[thick, ->] (14,0) -- (46.5,0);
  \foreach \v in {15,20,...,45} { \draw (\v,0.08) -- (\v,-0.08) node[below, font=\scriptsize] {\v}; }
  \node[font=\scriptsize\itshape] at (30,-0.68) {delivery time (minutes) --- one common scale};
  % parallel boxplots: Tony's (top) and Rosa's (bottom)
  \draw[thick] (18,1.55) -- (22,1.55);  \draw[thick] (18,1.4) -- (18,1.7);
  \draw[thick, fill=frost] (22,1.35) rectangle (30,1.75);
  \draw[very thick] (25,1.35) -- (25,1.75);
  \draw[thick] (30,1.55) -- (33,1.55);  \draw[thick] (33,1.4) -- (33,1.7);
  \fill[charcoal] (44,1.55) circle[radius=1mm];
  \draw[thick] (24,0.7) -- (31,0.7);    \draw[thick] (24,0.55) -- (24,0.85);
  \draw[thick, fill=goldbg] (31,0.5) rectangle (37,0.9);
  \draw[very thick] (35,0.5) -- (35,0.9);
  \draw[thick] (37,0.7) -- (39,0.7);    \draw[thick] (39,0.55) -- (39,0.85);
  \node[anchor=east, font=\footnotesize\bfseries] at (14,1.55) {Tony's};
  \node[anchor=east, font=\footnotesize\bfseries] at (14,0.7) {Rosa's};
  % parallel dot plots: Tony's (top) and Rosa's (bottom)
  \draw[slate] (15,2.45) -- (45,2.45);
  \draw[slate] (15,3.55) -- (45,3.55);
  \foreach \v/\k in {24/0, 29/0, 31/0, 33/0, 34/0, 35/0, 36/0, 36/1, 37/0, 38/0, 39/0} {
    \fill[goldacc] (\v,{2.45+0.13+0.24*\k}) circle[radius=1mm]; }
  \foreach \v/\k in {18/0, 21/0, 22/0, 24/0, 25/0, 25/1, 27/0, 28/0, 30/0, 33/0, 44/0} {
    \fill[cerulean] (\v,{3.55+0.13+0.24*\k}) circle[radius=1mm]; }
  \node[anchor=east, font=\footnotesize\bfseries] at (14,2.55) {Rosa's};
  \node[anchor=east, font=\footnotesize\bfseries] at (14,3.65) {Tony's};
\end{tikzpicture}\par\vspace{5pt}
Top pair: parallel \labelbox{2.6cm}{}\qquad Bottom pair: parallel \labelbox{2.6cm}{}\par}
\\
& \notesprompt{Read the displays --- every answer names the shop and the minutes.}
\begin{probgrid}{|Y|Y|}
\pcell{5}{Whose middle half of deliveries sits farther right on the common scale? What does that mean for a customer?}{2.0cm}{} & \pcell{6}{Find each shop's IQR from its box. Whose middle half is more spread out?}{2.0cm}{} \\ \hline
\end{probgrid}
\\
& \begin{probgrid}*{|Y|Y|}
\pcell{7}{Tony's $44$ is plotted as an outlier. Is Rosa's $24$-minute delivery one? Use the lower fence $Q_1 - 1.5 \cdot \text{IQR}$.}{2.0cm}{} & \pcell{8}{Mia draws Tony's boxplot on a $15$--$45$ scale and Rosa's on its own $20$--$40$ scale. Rosa's box now looks wider. Is Rosa's more spread out?}{2.4cm}{} \\ \hline
\end{probgrid}
\\ \hline
% ── ROW 3: comparing — shape, outliers, center, spread, in context ───────────
\mainidea[Shape, outliers, center, spread]{Compare in Context} &
A \textbf{comparison} compares all four features --- shape, outliers (unusual features), center,
and spread --- using comparative words (\emph{greater than, less than, about the same as}), in
context. \textbf{Pick the measures that fit the shape:} median and IQR when a distribution is
skewed or has an outlier; mean and standard deviation when both are roughly symmetric with no
outliers.
\\
& \textbf{9.} Fill in the table from rows 1 and 2.
\wordbank{skewed right \;\textbullet\; skewed left \;\textbullet\; 44 min \;\textbullet\; none \;\textbullet\; 25 min \;\textbullet\; 35 min \;\textbullet\; 8 min \;\textbullet\; 6 min}
{\renewcommand{\arraystretch}{1.5}%
\begin{tabularx}{\linewidth}{@{} X p{3.2cm} p{3.2cm} @{}}
\toprule
\textbf{Feature} & \textbf{Tony's} & \textbf{Rosa's} \\
\midrule
Shape & \blank{2.8cm} & \blank{2.8cm} \\
Outliers (unusual features) & \blank{2.8cm} & \blank{2.8cm} \\
Center (median) & \blank{2.8cm} & \blank{2.8cm} \\
Spread (IQR) & \blank{2.8cm} & \blank{2.8cm} \\
\bottomrule
\end{tabularx}}
\\
& \fcolorbox{cerulean}{frost}{\parbox{\dimexpr\linewidth-2\fboxsep-2\fboxrule\relax}{\raggedright
\textbf{The frame:} [Group A]'s [measure] ([value]) is \emph{greater than / less than / about
the same as} [Group B]'s ([value]), so [what that means in context].}}
\notesprompt{Use the frame --- one sentence each, in context.}
\begin{probgrid}{|Y|Y|}
\pcell{10}{Compare the centers.}{2.0cm}{} & \pcell{11}{Compare the spreads.}{2.0cm}{} \\ \hline
\end{probgrid}
\\
& \begin{probgrid}*{|Y|}
\pcell{12}{Compare the shapes and the unusual features.}{2.0cm}{} \\ \hline
\end{probgrid}
\\ \hline
% ── ROW 4: THE CRUX — same center, different spread ──────────────────────────
\mainidea[Same center,]{Different Spread} &
Back to the hook. Rosa's Pizza and Slice House each made $11$ deliveries, and both have a
median of $35$ minutes. \textbf{Your party needs the pizza within $40$ minutes.}
\\
& {\centering
\begin{tikzpicture}[x=0.21cm, y=1cm]
  \draw[thick, ->] (9,0) -- (62,0);
  \foreach \v in {10,15,...,60} { \draw (\v,0.08) -- (\v,-0.08); }
  \foreach \v in {10,20,...,60} { \node[below, font=\scriptsize] at (\v,-0.08) {\v}; }
  \node[font=\scriptsize\itshape] at (35,-0.68) {delivery time (minutes) --- one common scale};
  % both medians
  \draw[redacc, thick, densely dashed] (35,0) -- (35,2.2);
  \node[above, font=\scriptsize\bfseries\color{redacc}] at (35,2.2) {both medians: $35$ min};
  % Slice House (bottom) and Rosa's (top)
  \draw[slate] (10,0.5) -- (60,0.5);
  \draw[slate] (10,1.5) -- (60,1.5);
  \foreach \v in {15, 20, 24, 28, 33, 35, 38, 44, 47, 50, 55} {
    \fill[plumacc] (\v,0.63) circle[radius=1mm]; }
  \foreach \v/\k in {24/0, 29/0, 31/0, 33/0, 34/0, 35/0, 36/0, 36/1, 37/0, 38/0, 39/0} {
    \fill[goldacc] (\v,{1.5+0.13+0.24*\k}) circle[radius=1mm]; }
  \node[anchor=south west, font=\footnotesize\bfseries] at (10,0.78) {Slice House};
  \node[anchor=south west, font=\footnotesize\bfseries] at (10,1.55) {Rosa's};
\end{tikzpicture}\par}
\\
& \fcolorbox{cerulean}{frost}{\parbox{\dimexpr\linewidth-2\fboxsep-2\fboxrule\relax}{\centering
\textbf{The same center does not make the same distribution. When the centers match, the
spread decides: the group with less spread is more consistent.}}}

\vspace{3pt}
A list of two separate descriptions --- or a comparison of centers alone --- is \emph{not} a
comparison.\par
\textbf{In your own words:} why can two shops with the same median not be equally good for
your party?
\writespace{1.4cm}{}
\\
& \notesprompt{Now settle the hook vote.}
\begin{probgrid}{|Y|Y|}
\pcell{13}{Find each shop's median and IQR. (Slice House: the median of each half of its $11$ times.)}{2.0cm}{} & \pcell{14}{\textbf{Vote again:} Rosa's \; Slice House \; doesn't matter. How many of each shop's $11$ deliveries came within $40$ minutes? Which shop?}{2.4cm}{} \\ \hline
\end{probgrid}
\\
& \begin{probgrid}*{|Y|Y|}
\pcell{15}{Another night you only care about your best chance at a delivery \emph{under $25$ minutes}. Now which shop?}{2.0cm}{} & \pcell{16}{Dev writes: \emph{``Rosa's: median 35, IQR 6. Slice House: median 35, IQR 23.''} Is that a comparison? Rewrite it, in context.}{2.6cm}{} \\ \hline
\end{probgrid}
\\ \hline
% ── ROW 5: GUIDED PRACTICE — we do, together, a fresh context ────────────────
\mainidea[Guided practice]{Paper Airplanes} &
The design club throws two paper-airplane designs $10$ times each and measures every flight in
feet. Below are the calculator's \emph{1-Var Stats} screens and a parallel dot plot. Sx is the
\emph{sample} standard deviation --- it divides by $n - 1$.
\\
& {\centering
\begin{tikzpicture}
  \foreach \x/\des/\sd/\lo/\qa/\qc/\hi in {0/A/2.58/26/28/32/34, 6/B/7.26/18/25/35/42} {
    \fill[charcoal, rounded corners=3pt] (\x,0) rectangle ({\x+5.2},3.2);
    \fill[frost] ({\x+0.15},0.15) rectangle ({\x+5.05},3.05);
    \node[anchor=north west, align=left, font=\ttfamily\footnotesize] at ({\x+0.35},2.95)
      {\textbf{1-Var Stats: Design \des}\\ $\bar{x}$=30\\ Sx=\sd\\ n=10\\ minX=\lo \quad Q1=\qa\\ Med=30 \quad Q3=\qc\\ maxX=\hi};
  }
\end{tikzpicture}\par}
\\
& {\centering
\begin{tikzpicture}[x=0.3cm, y=1cm]
  \draw[thick, ->] (14,0) -- (46.5,0);
  \foreach \v in {15,20,...,45} { \draw (\v,0.08) -- (\v,-0.08) node[below, font=\scriptsize] {\v}; }
  \node[font=\scriptsize\itshape] at (30,-0.68) {flight distance (feet) --- one common scale};
  \draw[slate] (15,0.5) -- (45,0.5);
  \draw[slate] (15,1.6) -- (45,1.6);
  \foreach \v/\k in {18/0, 22/0, 25/0, 28/0, 30/0, 30/1, 32/0, 35/0, 38/0, 42/0} {
    \fill[plumacc] (\v,{0.5+0.13+0.24*\k}) circle[radius=1mm]; }
  \foreach \v/\k in {26/0, 27/0, 28/0, 29/0, 30/0, 30/1, 31/0, 32/0, 33/0, 34/0} {
    \fill[cerulean] (\v,{1.6+0.13+0.24*\k}) circle[radius=1mm]; }
  \node[anchor=east, font=\footnotesize\bfseries] at (14,0.6) {Design B};
  \node[anchor=east, font=\footnotesize\bfseries] at (14,1.7) {Design A};
\end{tikzpicture}\par}
\\
& \notesprompt{We work these together.}
\begin{probgrid}{|Y|Y|}
\pcell{17}{Shape and outliers: are both distributions roughly symmetric with no outliers? So which pair of measures fits?}{2.0cm}{} & \pcell{18}{Compare the centers, then the spreads, with comparative words and units.}{2.2cm}{} \\ \hline
\end{probgrid}
\\
& \begin{probgrid}*{|Y|Y|}
\pcell{19}{The contest target is a landing zone from $27$ to $33$ feet. Kai says, \emph{``Same mean --- it doesn't matter which design we fly.''} Count each design's throws in the zone. Is Kai right?}{2.6cm}{} & \pcell{20}{Write the full comparison --- shape, outliers, center, spread --- in context, in one or two sentences.}{2.6cm}{} \\ \hline
\end{probgrid}
\\ \hline
\end{guidednotes}

% ── THE NOTES END AT GUIDED PRACTICE ─────────────────────────────────────────
% There is deliberately no "you do" here: ../homework/ is the individual practice,
% assigned at Close & Assign and done outside class.

\end{document}
